\clearpage\pagestyle{otMain}\fontsize{11}{15}\selectfont\corrigesChapitre{3}{Multiples, diviseurs et nombres premiers}
\begin{blocLivre}{cor-3-1}\begin{corrige}{1}{Tester des diviseurs}\demarche{Utiliser les critères puis contrôler les quotients.}
2340 est pair et se termine par 0 : 2, 5 et 10 conviennent. La somme de ses chiffres vaut 9 : 3 et 9 conviennent aussi. Les cinq nombres proposés sont donc des diviseurs.
\end{corrige}\end{blocLivre}
\begin{blocLivre}{cor-3-2}\begin{corrige}{2}{Premier ou composé ?}\demarche{Appliquer la définition et fournir des facteurs quand c’est possible.}
$\res{2\text{ et }17}$ sont premiers. Le nombre 1 n’a qu’un diviseur positif ; $21=3\times7$ et $49=7^2$ sont composés. Pour 17, les seuls diviseurs premiers à tester sont 2 et 3, car $5^2=25>17$ ; ils ne divisent pas 17.
\end{corrige}\end{blocLivre}
\begin{blocLivre}{cor-3-3}\begin{corrige}{3}{Décomposer un entier}\demarche{Poursuivre la décomposition jusqu’aux facteurs premiers.}
$420=42\times10=(2\times3\times7)(2\times5)=\res{2^2\times3\times5\times7}$. Tous les facteurs sont premiers ; leur produit redonne 420.
\end{corrige}\end{blocLivre}
\begin{blocLivre}{cor-3-4}\begin{corrige}{4}{Une fraction irréductible}\demarche{Simplifier tous les facteurs premiers communs.}
Les décompositions sont $198=2\times3^2\times11$ et $252=2^2\times3^2\times7$. En simplifiant $2\times3^2$, on obtient $\res{11/14}$. Les diviseurs de 11 ne divisent pas 14.
\end{corrige}\end{blocLivre}
\begin{blocLivre}{cor-3-5}\begin{corrige}{5}{Écrire des multiples}\demarche{Encadrer les multiplicateurs entiers.}
On part de $14\times5=70$ puis on ajoute 14 : $\res{70,84,98,112,126}$. Le suivant vaut 140 et dépasse la borne supérieure.
\end{corrige}\end{blocLivre}
\begin{blocLivre}{cor-3-6}\begin{corrige}{6}{Tous les diviseurs}\demarche{Chercher des paires de facteurs sans doublons.}
On utilise les paires $1\times72$, $2\times36$, $3\times24$, $4\times18$, $6\times12$ et $8\times9$. On obtient $\res{1,2,3,4,6,8,9,12,18,24,36,72}$. Au-delà de 8, chaque diviseur a déjà son partenaire dans une paire, car $9\times9>72$.
\end{corrige}\end{blocLivre}
\begin{blocLivre}{cor-3-7}\begin{corrige}{7}{Un chiffre inconnu}\demarche{Combiner les critères, puis vérifier les candidats.}
La divisibilité par 2 impose $a\in\{0,2,4,6,8\}$. La somme des chiffres vaut $9+a$ ; elle est divisible par 3 pour $a=0,3,6,9$. Les chiffres communs aux deux listes donnent $\res{a=0\text{ ou }6}$.
\end{corrige}\end{blocLivre}
\begin{blocLivre}{cor-3-8}\begin{corrige}{8}{Une grande fraction}\demarche{Soustraire les exposants seulement pour les facteurs de même base.}
On simplifie $2^2$ et $3^2$. Il reste $(2^2\times7)/(3\times5)=\res{28/15}$. Aucun facteur 2 ou 7 ne peut encore se simplifier avec 3 ou 5.
\end{corrige}\end{blocLivre}
\begin{blocLivre}{cor-3-9}\begin{corrige}{9}{Un carré parfait ?}\demarche{Observer la parité de chaque exposant premier.}
Un carré a des exposants pairs dans sa décomposition. L’exposant 3 du facteur 3 est impair : le nombre n’est pas un carré. Multiplier par $\res3$ donne $2^4\times3^4=(2^2\times3^2)^2=\res{36^2}$. Un facteur 3 est nécessaire.
\end{corrige}\end{blocLivre}
\begin{blocLivre}{cor-3-10}\begin{corrige}{10}{Un calcul par tableur}\demarche{Relier le reste nul à la définition d’un diviseur.}
La fonction MOD calcule le reste de la division par la valeur de A2. Un reste nul signifie que cet entier divise 84. La liste cherchée est $\res{1,2,3,4,6,7,12,14}$. On ne trouve pas les diviseurs supérieurs à 20 avec cette seule colonne.
\end{corrige}\end{blocLivre}
\begin{blocLivre}{cor-3-11}\begin{corrige}{11}{Préparer des lots}\demarche{Traduire le partage en diviseurs communs.}
Le nombre de lots doit diviser 156 et 204. Or $156=2^2\times3\times13$ et $204=2^2\times3\times17$ : le plus grand diviseur commun vaut 12. On peut former $\res{12\text{ lots}}$, contenant chacun $\res{13\text{ rouges et }17\text{ bleues}}$. Aucun diviseur commun n’est plus grand.
\end{corrige}\end{blocLivre}
\begin{blocLivre}{cor-3-12}\begin{corrige}{12}{Des cigales périodiques}\demarche{Chercher le premier multiple commun aux deux périodes.}
Le nombre d’années écoulées doit être un multiple commun de 13 et de 17. Ces deux nombres sont premiers et distincts : le plus petit multiple commun positif est $13\times17=221$. La prochaine coïncidence aura lieu en $\res{2245}$. Des périodes premières entre elles rendent les rencontres rares.
\end{corrige}\end{blocLivre}
\begin{blocLivre}{cor-3-13}\begin{corrige}{13}{Daller un rectangle}\demarche{Utiliser un diviseur commun pour les deux dimensions.}
Le côté doit diviser 180 et 252. Leur plus grand diviseur commun vaut 36. On place $252/36=7$ carrés sur la longueur et $180/36=5$ sur la largeur : $\res{35\text{ carrés de côté }36\text{ cm}}$. Le contrôle par les aires donne $35\times36^2=180\times252$.
\label{fig-cor-3-13}\figureGeo[corrige]{dallage}
\end{corrige}\end{blocLivre}
\begin{blocLivre}{cor-3-14}\begin{corrige}{14}{Un stock à retrouver}\demarche{Tester la compatibilité des contraintes avant de chercher un nombre.}
Un tel nombre doit être multiple de $2^3$, de $3^2$ et de 5, donc de $8\times9\times5=360$. Les multiples voisins sont 360 et 720. Aucun n’est entre 500 et 600 : $\res{\text{aucune solution}}$. Un énoncé peut conduire légitimement à l’impossibilité.
\end{corrige}\end{blocLivre}
\begin{blocLivre}{cor-3-15}\begin{corrige}{15}{Des engrenages}\demarche{Compter les dents au contact plutôt que confondre les nombres de tours.}
Le nombre de dents passées au contact doit être multiple de 24 et de 40. Le premier multiple commun est 120. La petite roue fait $120/24=\res5$ tours et la grande $120/40=\res3$ tours, en sens opposés. Le contact initial revient après $\res{120\text{ passages}}$.
\label{fig-cor-3-15}\figureGeo[corrige]{engrenages}
\end{corrige}\end{blocLivre}
